\section{Complex references}\label{sec:complex-references}

The activities of the renormalization group must be analytic in the background on a complex domain (blueprint, Appendix~F, domain D3$'''$). So the step evaluates the operators of Part~IX, the interior estimates and the weakened background map at \emph{complex} references. Sections~\ref{sec:interior}--\ref{sec:weak-background} were written for real backgrounds. This section records where reality was used, and shows that every statement needed later extends to small complex perturbations. The multiscale averaged backgrounds of Part~IV need a separate argument, given in Section~\ref{sec:AB-proof} (Proposition~\ref{prop:AB}). The results of Part~IX are extended in Section~\ref{sec:complex-ext}.

\paragraph{Complex references.} Let $U_r$ be a real $(M,\eps)$-regular background. A \emph{complex reference} of size $\delta_c$ is $U=e^{\ii\eta A'}U_r$, with $A'$ a $\mathfrak g^c$-valued field such that, in the local real gauges of Definition~\ref{def:regular} for $U_r$,
\[
|A'|+|\nabla A'|+\sup_y[\nabla\tilde A'_y]_{\alpha,\Lam_1(c_y)}<\delta_c .
\]
\begin{itemize}
\item \emph{Relation to Proposition~\ref{prop:reg-complex}.} That proposition writes $U\,e^{\ii\eta\,\delta\!A}$ and uses ordinary derivatives. The two forms are equivalent, with $\delta\!A=R(U_r)^{-1}A'$. The two sizes are comparable: there is $\eps_{\rm cc}>0$, depending only on $d$, $\alpha$ and $G$, such that for $\eps\le\eps_{\rm cc}$ the $\delta$ of Proposition~\ref{prop:reg-complex} and the size of $A'$ above bound each other up to the factor $C_{\rm cmp}=4$. Indeed, in the gauge of Definition~\ref{def:regular} on a cube $\hat\Lam$, where $U_r=e^{\ii\eta A_r}$, the two fields differ by the factor $R(e^{\ii\eta A_r})$, which is within $4\eta\eps$ of $1$, with difference quotients and H\"older quotients of size $C\eta\eps$. Covariant and flat gradients differ by $C\eps|A'|$, by \eqref{eq:cov-vs-flat}, and the change between the gauges of the blocks and of the cubes costs at most $C\eps$ times the $C^1$ norm, as in the proof of Theorem~\ref{thm:ID}(b). Finally, two points at $\ell^\infty$-distance less than $1$ lie in a common cube $\Lam_1(c_y)$, and pairs at distance at least $1$ have H\"older quotients at most twice the sup of the gradient.
\item \emph{The local gauge.} In a local gauge of $U_r$ we have $U=e^{\ii\eta A_0}$, with $A_0=A'\oplus A_r$ combined bondwise by the Baker--Campbell--Hausdorff formula. Then $A_0$ is complex and satisfies (B) with parameter $C(\eps+\delta_c)$.
\item \emph{Potentials.} Local potentials at a complex reference are taken in $\cW^{\rm an}_{\rm loc}$ (Section~\ref{sec:reg-gauge}), and satisfy the bounds of $\cW_{\rm loc}(\eps+C\delta_c)$ in the real gauges of $U_r$. The non-local potential $W_2^X$ satisfies the hypotheses of Proposition~\ref{prop:reg-complex} for existence, $\Rea$-coercivity and decay; its comparison of domains is Proposition~\ref{prop:dirichlet-hessian}(d). Only its $L^2$ statements are used (Section~\ref{sec:localization}). In Section~\ref{sec:weak-background}, $W=0$.
\item \emph{The domain D3$'''$.} It only needs complex perturbations near a polymer. These are extended by $U_r$ outside a neighbourhood, with a smooth cutoff at unit scale, which gives a complex reference with $\delta_c\le Ch^{1-2\epsilon}$. In D3$'''$ the real part has only curvature and current at most $h^{1-\epsilon/2}$. Its $(M,\eps)$-regularity comes from Corollary~\ref{cor:regularity}, with constants depending on $L$ and $M$. Its current is at most $\eps$ for $g_*$ small; Section~\ref{sec:step-quantities} adds $|J_{U_r}|\le\eps$ to the definition of a complex reference, for the coercivity of the Hessian.
\end{itemize}

\subsection{Averaged backgrounds at complex references}\label{sec:AB-proof}

Proposition~\ref{prop:Q-expansion} and Proposition~\ref{prop:linearized-averages} assume a real $U_0$, with real averaged backgrounds $W_j\in\mathfrak U_j(2a)$ (Proposition~\ref{prop:avg-regular}). Lemma~\ref{lem:one-step-expansion} uses this through:
\begin{itemize}
\item the logarithms of Lemma~\ref{lem:one-step-basic}(a);
\item the curvature bound on $H_x$, $E$;
\item Lemma~\ref{lem:avg-local-gauge}.
\end{itemize}
The underlying lemmas are proved only for real configurations: Lemma~\ref{lem:log-estimates}(b), with the bi-invariant metric on $G$; Lemma~\ref{lem:stokes}; and Lemma~\ref{lem:one-step-regularity}, which uses a real gauge and gauge invariance of $|F|$. These results feed the bounds $|\bQ A(c)|\le C\|A\|$ and $|Q'\lambda|\le C\|\lambda\|$, and the bounded transposes, used in Theorem~\ref{thm:ID} and in Part~IX. A naive per-scale bound on complex transports would compound to $2^{j}$ over the scales.

We therefore work relative to the real reference $U_r$. Two observations make this work.
\begin{itemize}
\item The averaged backgrounds at $U$ are gauge transforms of $X_jW_j(U_r)$, with $X_j=\mathcal Q_j(U')$, by explicit complex gauge transformations $h_j$ (Lemma~\ref{lem:AB-rep}). This is proved directly from the one-step definitions. The axial gauge, Lemma~\ref{lem:axial}(b),(c) and Theorem~\ref{thm:correspondence}(a),(b),(e),(f) are not used at complex configurations.
\item At the points of $\T^{(j)}$, $h_j$ has size $C\delta_j|A'|$, not $C|A'|$. It is a product of the contour variables $\tau_{X_i}$, $i<j$. Each $\tau_{X_i}$ is the exponential of an average of logarithms of relative holonomies along contours of length at most $d\delta_{i+1}$, and $\sum_{i<j}\delta_{i+1}\le2\delta_j$. The scale-$j$ gradient of $h_j$ carries the factor $\delta_j^{-1}$, so it is bounded by the sup norm. Hence the relative fields are $O(|A'|)$ uniformly in $j$, and the losses per scale are summable.
\end{itemize}
Only the sup norm of $A'$ is used, and the smallness of $\delta_c$ depends on $d$, $L$ and the group $G\subset U(\mathsf N)$, including the radius $c_G$ of its logarithm chart, but not on $M$.

\paragraph{Setting.} Let $U_r$ be real and $(M,\eps)$-regular with $\eps\le\eps_0$. By the proof of Corollary~\ref{cor:parts-apply}, $U_r\in\mathfrak U_0(a)$ with $a\le C\eps$, where $C$ does not depend on $M$ (Lemma~\ref{lem:curvature}). The proofs of Sections~\ref{sec:averaging} and~\ref{sec:adapted} do not use $M$. So $c_3$ and the constants $C,c$ of this subsection depend only on $d$, $L$ and $G$. Let $c'_3\le c_3$ be such a constant, decreased below where needed, and assume $a\le c'_3$. Put $W^r_j=W_j(U_r)=U_r^{(j)}$. These are real and lie in $\mathfrak U_j(2a)$ by Proposition~\ref{prop:avg-regular}.

Let $U=U'U_r$ (bondwise), with $U'=e^{\ii\eta A'}$ and $A'$ a $\mathfrak g^c$-valued field. Under a real gauge transformation $u$, $U_r\mapsto U_r^u$ and $A'(b)\mapsto R(u(b_-))A'(b)$. So $|A'|_\infty$ is gauge invariant, and the condition $|A'|<\delta_c$ in the local gauges of $U_r$ is the condition $|A'|_\infty<\delta_c$. The derivative bounds on $A'$ in the definition of complex references are not used in this subsection. Let $W_j(U)=U^{(j)}=\mathfrak a^{(j)}(U)$, whenever defined.

We use the gauge-adapted averaging of Section~\ref{sec:adapted} relative to the \emph{real} reference $U_r$, so there $W_j=W^r_j$. Applied to the relative configuration $U'$ it gives
\[
X_j=\mathcal Q_j(U')=e^{\ii\delta_jB_j},\qquad \tau_i=\tau_{X_i}\ \ (\text{on }\T^{(i+1)}),\qquad h_j(y)=\tau_0(y)\cdots\tau_{j-1}(y)\ \ (y\in\T^{(j)}),
\]
with $h_0=1$. By Theorem~\ref{thm:correspondence}(c), which holds for complex relative configurations (remark after that theorem), $h_j=\mathsf w_j(U')^{-1}$ at the points of $\T^{(j)}$. As in Definition~\ref{def:adapted} we use $\tilde{\mathfrak a}_j(X)=\overline X^{(\tau_X^{-1})}$. By Lemma~\ref{lem:AB-cov} below, this coincides with $\overline{X^{(g_X)}}$ on the domain of Lemma~\ref{lem:one-step-expansion}.

Norms of linear maps of $\mathfrak g^c$, and of matrices acting on $\C^{\mathsf N}$, are operator norms, written $\|\cdot\|$. The constants $C_{\rm AB}=14d+3$, $C_{\rm loop}$ and $\kappa_{\rm AB}=12dC_{\rm AB}$ depend only on $d$ and $\mathsf N$. Recall $\delta_j=L^{j-k}$, so that
\begin{equation}\label{eq:AB-scales}
\sum_{i\le k}\delta_i\le2,\qquad\sum_{i<j}\delta_{i+1}=\delta_j\sum_{m<j}L^{-m}\le2\delta_j .
\end{equation}

\begin{lemma}[Covariance under complex gauge transformations]\label{lem:AB-cov}
Let $j\ge1$ and let $Y$ be a configuration on $\T^{(j-1)}$ with $|Y(\Gamma^j_{c,x})Y([c])^{-1}-1|\le\frac12c_G$ for all bonds $c$ of $\T^{(j)}$ and all $x\in B_j(c_-)$. Let $g:\T^{(j-1)}\to G^c$ with $\|g(z)\|\|g(z)^{-1}\|\le2$ for all $z$. Then $\mathfrak a_j(Y)$ and $\mathfrak a_j(Y^g)$ are defined, and $\mathfrak a_j(Y^g)=\mathfrak a_j(Y)^{g|_{\T^{(j)}}}$.

Consequently, let $X$ be a relative configuration at scale $j-1$ such that $XW_{j-1}$ satisfies the hypothesis on $Y$, for some background $W_{j-1}$ with $W_j=\mathfrak a_j(W_{j-1})$. Then $\overline{X^{(g)}}=\overline X^{(g|_{\T^{(j)}})}$.
\end{lemma}

\begin{proof}
The loops of $Y^g$ are $g(c_-)hg(c_-)^{-1}$, with $h$ the corresponding loop of $Y$, and $|ghg^{-1}-1|\le\|g\|\|g^{-1}\||h-1|\le c_G$. Since $(ghg^{-1}-1)^n=g(h-1)^ng^{-1}$, the power series give $\ell(ghg^{-1})=g\,\ell(h)\,g^{-1}$, and $\exp(\ii gZg^{-1})=g\,e^{\ii Z}g^{-1}$. Insert this and $Y^g([c])=g(c_-)Y([c])g(c_+)^{-1}$ into Definition~\ref{def:one-step}. For relative configurations use $X^{(g)}W_{j-1}=(XW_{j-1})^g$.
\end{proof}

This makes the remark after Lemma~\ref{lem:one-step-basic} quantitative, without analytic continuation.

\begin{lemma}[Representation of the averaged backgrounds]\label{lem:AB-rep}
There is $\delta_{AB}>0$, depending on $d,L,G$, such that for $a\le c'_3$ and $|A'|_\infty\le\delta_{AB}$:
\begin{enumerate}[label=(\alph*)]
\item $|\ell(\tau_i)|\le3d\,\delta_{i+1}|A'|_\infty$ on $\T^{(i+1)}$ for $i<k$. The function $\lambda_j=\ell(h_j)$ on $\T^{(j)}$ is defined and satisfies $|\lambda_j|\le7d\,\delta_j|A'|_\infty$ for $j\le k$.
\item For every $j\le k$, $W_j(U)$ is defined and
\begin{equation}\label{eq:AB-rep}
W_j(U)=\big(X_jW^r_j\big)^{h_j},\qquad\text{i.e.}\quad W_j(U)(b)=h_j(b_-)\,X_j(b)\,W^r_j(b)\,h_j(b_+)^{-1}.
\end{equation}
Equivalently, $U'^{[j]}=\mathcal Q_j(U')^{(h_j)}=\mathcal Q_j(U')^{(\mathsf w_j(U')^{-1})}$.
\end{enumerate}
\end{lemma}

\begin{proof}
(a) Proposition~\ref{prop:Q-expansion}(a) applies to the real reference $U_r$ and the complex field $A'$, and gives $|B_i|_\infty\le2|A'|_\infty$. By Lemma~\ref{lem:one-step-expansion}, $\ell(\tau_i)=\delta_{i+1}\big(\gamma_{i+1}B_i+c^\tau_{i+1}(B_i)\big)$ with $|\gamma_{i+1}B_i|\le2d|A'|$ and $|c^\tau_{i+1}(B_i)|\le4C\delta_{i+1}|A'|^2$. This gives the first bound for $|A'|\le d/(4C)$. By \eqref{eq:AB-scales}, $s=\sum_{i<j}|\ell(\tau_i)|\le6d\delta_j|A'|$. Lemma~\ref{lem:log-estimates}(c) gives $|\lambda_j|\le s+Cs^2\le7d\delta_j|A'|$ for $|A'|$ small.

(b) We use induction on $j$. For $j=0$ we have $W_0(U)=U=U'U_r=X_0W^r_0$ and $h_0=1$. Let $j\ge1$, assume \eqref{eq:AB-rep} at $j-1$, and put $Y=X_{j-1}W^r_{j-1}$. With $h^r_x=W^r_{j-1}(\Gamma^j_{c,x})W^r_{j-1}([c])^{-1}$,
\[
Y(\Gamma^j_{c,x})Y([c])^{-1}=(X_{j-1})_{W^r}(\Gamma^j_{c,x})\;h^r_x\;(X_{j-1})_{W^r}([c])^{-1}.
\]
The relative transports are products of at most $(2d+2)L$ factors $e^{\ii\delta_{j-1}R(W^r(\cdot))B_{j-1}(\cdot)}$, and $R(W^r)$ is unitary. Moreover $|\ell(h^r_x)|\le2da\delta_j^2$, by Lemma~\ref{lem:one-step-basic}(a) and Proposition~\ref{prop:avg-regular}. By Lemma~\ref{lem:log-estimates}(c) the loop has logarithm at most $2\big(4(d+1)\delta_j|A'|+2da\delta_j^2\big)$. So it is within $\frac14c_G$ of $1$, after decreasing $c'_3$ and $\delta_{AB}$. By (a), $\|h_{j-1}\|\|h_{j-1}^{-1}\|\le e^{2|\lambda_{j-1}|}\le2$. So the loops of $W_{j-1}(U)=Y^{h_{j-1}}$, which are conjugates of those of $Y$ by $h_{j-1}(c_-)$, are within $\frac12c_G$ of $1$; this is used in Lemma~\ref{lem:AB-bounds}(b). By the induction hypothesis $W_{j-1}(U)=Y^{h_{j-1}}$, so Lemma~\ref{lem:AB-cov} shows that $W_j(U)=\mathfrak a_j(W_{j-1}(U))$ is defined and
\[
W_j(U)=\mathfrak a_j(Y)^{h_{j-1}|_{\T^{(j)}}} .
\]
By the definitions of the relative average and of $\tilde{\mathfrak a}_j$,
\[
\mathfrak a_j(Y)=\overline{X_{j-1}}\,W^r_j,\qquad \overline{X_{j-1}}=\tilde{\mathfrak a}_j(X_{j-1})^{(\tau_{j-1})}=X_j^{(\tau_{j-1})},
\]
i.e.\ $\mathfrak a_j(Y)=(X_jW^r_j)^{\tau_{j-1}}$. Since $(Z^g)^{g'}=Z^{g'g}$ and $h_{j-1}\tau_{j-1}=h_j$ on $\T^{(j)}$, \eqref{eq:AB-rep} follows at $j$. The last statement is \eqref{eq:AB-rep} rewritten with the relative gauge action, together with $h_j=\mathsf w_j(U')^{-1}$ on $\T^{(j)}$.
\end{proof}

\begin{lemma}[Averaged backgrounds at complex references]\label{lem:AB-bounds}
Possibly after decreasing $\delta_{AB}$ and $c'_3$ (depending on $d,L,G$), the following holds for $a\le c'_3$, $|A'|_\infty\le\delta_c\le\delta_{AB}$ and $0\le j\le k$.
\begin{enumerate}[label=(\alph*)]
\item $W_j(U)=e^{\ii\delta_j\tilde B_j}\,W^r_j$ bondwise, where
\begin{equation}\label{eq:AB-field}
\tilde B_j=B_j-D_{W^r_j}\lambda_j+\beta_j,\qquad|\beta_j|_\infty\le C\delta_j|A'|^2_\infty,\qquad|\tilde B_j|_\infty\le C_{\rm AB}|A'|_\infty .
\end{equation}
Moreover $\big|\tilde B_j-\bQ_j(U_r)A'+D_{W^r_j}\boldsymbol\Gamma_jA'\big|_\infty\le C|A'|_\infty^2$, where $\boldsymbol\Gamma_jA'=\sum_{i<j}\delta_{i+1}\gamma_{i+1}\bQ_i(U_r)A'$ is evaluated at the points of $\T^{(j)}$.
\item $\tilde B_j$ is holomorphic in $A'$ on $\{|A'|_\infty<\delta_{AB}\}$, and real for real $A'$. Under a real gauge transformation $u$ of $\Tf$ ($U_r\mapsto U_r^u$, $A'(b)\mapsto R(u(b_-))A'(b)$) we have $\tilde B_j(b)\mapsto R(u(b_-))\tilde B_j(b)$. In particular (a) holds in every real gauge, e.g.\ in the local gauges of $U_r$.
\item Let $\gamma$ be a contour of $n$ bonds of $\T^{(j)}$ with $s=C_{\rm AB}n\delta_j|A'|_\infty\le s_0$, where $s_0$ depends only on $\mathsf N$. Then $W_j(U)(\gamma)=e^{\ii Z_\gamma}W^r_j(\gamma)$ with $|Z_\gamma|\le2s$. (Here $Z_\gamma$ is the Baker--Campbell--Hausdorff series of $\mathfrak g^c$-elements of Lemma~\ref{lem:log-estimates}(c), which lies in $\mathfrak g^c$; so no radius $c_G$ enters.) Hence
\[
\|R(W_j(U)(\gamma))^{\pm1}\|\le e^{4s},\qquad\|R(W_j(U)(\gamma))-R(W^r_j(\gamma))\|\le e^{4s}-1 .
\]
Applied at index $j-1$ to the contours of at most $3dL$ bonds of $\T^{(j-1)}$ used in the step $j-1\to j$, where $4s\le12dLC_{\rm AB}\delta_{j-1}\delta_c=\kappa_{\rm AB}\delta_j\delta_c$, this gives these bounds with $e^{4s}$ replaced by $e^{\kappa_{\rm AB}\delta_j\delta_c}$.
\item For $j\ge1$, a bond $c$ of $\T^{(j)}$ and $x\in B_j(c_-)$, let $E(U)=L^{-d}\sum_xH_x(U)$ and
\[
H_x(U)=\ell\big(W_{j-1}(U)(\Gamma^j_{c,x})\,W_{j-1}(U)([c])^{-1}\big).
\]
Then
\begin{gather*}
|H_x(U)-H_x(U_r)|,\ \ |E(U)-E(U_r)|\le C_{\rm loop}\delta_j|A'|_\infty,\\
|H_x(U)|\le2da\delta_j^2+C_{\rm loop}\delta_j|A'|_\infty,\qquad|E(U)|\le2da\delta_j^2+C_{\rm loop}\delta_j|A'|_\infty .
\end{gather*}
\end{enumerate}
\end{lemma}

\begin{proof}
(a) By Lemma~\ref{lem:AB-rep}, $V_j=W_j(U)(W^r_j)^{-1}=X_j^{(h_j)}$, i.e.
\[
V_j(b)=e^{\ii\lambda_j(b_-)}\,e^{\ii\delta_jB_j(b)}\,e^{-\ii R(W^r_j(b))\lambda_j(b_+)} .
\]
The three exponents have total size $s\le(14d+2)\delta_j|A'|\le\frac14$, by Lemma~\ref{lem:AB-rep}(a) and Proposition~\ref{prop:Q-expansion}(a); here we used that $R(W^r_j)$ is unitary. By Lemma~\ref{lem:log-estimates}(c), $\ell(V_j(b))=\lambda_j(b_-)+\delta_jB_j(b)-R(W^r_j(b))\lambda_j(b_+)+\delta_j\beta_j(b)$ with $\delta_j|\beta_j|\le Cs^2$. This is \eqref{eq:AB-field}, with $D_{W^r_j}\lambda(b)=\delta_j^{-1}\big(R(W^r_j(b))\lambda(b_+)-\lambda(b_-)\big)$ as in Section~\ref{sec:adapted}. By Lemma~\ref{lem:AB-rep}(a), $|D_{W^r_j}\lambda_j|\le\delta_j^{-1}\big(|\lambda_j(b_+)|+|\lambda_j(b_-)|\big)\le14d|A'|$. Hence $|\tilde B_j|\le(2+14d+C\delta_j|A'|)|A'|\le C_{\rm AB}|A'|$.

For the linear part, use $B_j=\bQ_jA'+C_j(A')$ with $|C_j(A')|\le C|A'|^2$ (Proposition~\ref{prop:Q-expansion}(b)). Further, $\lambda_j=\sum_{i<j}\ell(\tau_i)+O(\delta_j^2|A'|^2)$ and $\ell(\tau_i)=\delta_{i+1}\gamma_{i+1}\bQ_iA'+O(\delta_{i+1}|A'|^2)$. By \eqref{eq:AB-scales}, $|\lambda_j-\boldsymbol\Gamma_jA'|\le C\delta_j|A'|^2$, and $D_{W^r_j}$ costs at most $2\delta_j^{-1}$.

(b) $B_j$ and $\tau_i$ are holomorphic in $A'$ by Proposition~\ref{prop:Q-expansion} and Lemma~\ref{lem:one-step-expansion}. Then $\lambda_j$ and $\tilde B_j$ are holomorphic, since they are obtained from these by $\ell$, products and Lemma~\ref{lem:log-estimates}(c). For real $A'$ all configurations involved are real. For covariance, $U^u=e^{\ii\eta R(u_-)A'}U_r^u$. Lemma~\ref{lem:AB-cov} with real $g$, applied at each scale (its hypothesis holds by the proof of Lemma~\ref{lem:AB-rep}(b)), gives $W_j(U^u)=W_j(U)^{u|_{\T^{(j)}}}$, and likewise $W_j(U_r^u)=(W^r_j)^{u|_{\T^{(j)}}}$. Hence $V_j(b)\mapsto u(b_-)V_j(b)u(b_-)^{-1}$.

(c) $W_j(U)(\gamma)=(V_j)_{W^r_j}(\gamma)\,W^r_j(\gamma)$, and $(V_j)_{W^r_j}(\gamma)=\prod_iR(W^r_j(\gamma_{<i}))\big(V_j(b_i)\big)$. For negatively oriented bonds, $V_j(-b)=R(W^r_j(b))^{-1}(V_j(b)^{-1})$. Since $R(W^r_j)$ is unitary, every factor is $e^{\ii\delta_jZ_i}$ with $|Z_i|\le C_{\rm AB}|A'|$. Lemma~\ref{lem:log-estimates}(c) gives $|Z_\gamma|\le s+Cs^2\le2s$ for $s\le s_0=\min(\frac14,C^{-1})$. By \eqref{eq:Ad-bound}, $\|R(e^{\pm\ii Z})\|\le e^{2|Z|}$ and $\|R(e^{\ii Z})-1\|\le e^{2|Z|}-1$, and $R(W^r_j(\gamma))$ is unitary. The last statement is the case $n\le3dL$ at index $j-1$, with $\delta_{j-1}=\delta_j/L$; it needs $3dC_{\rm AB}\delta_j\delta_c\le s_0$, which holds after decreasing $\delta_{AB}$.

(d) We have
\[
W_{j-1}(U)(\Gamma^j_{c,x})\,W_{j-1}(U)([c])^{-1}=e^{\ii Z_1}\,e^{\ii H}\,e^{-\ii Z_2},
\]
with $H=H_x(U_r)$, $e^{\ii Z_1}=(V_{j-1})_{W^r}(\Gamma^j_{c,x})$ and $e^{\ii Z_2}=(V_{j-1})_{W^r}([c])$. By (c), $|Z_1|+|Z_2|\le4(d+1)C_{\rm AB}\delta_j|A'|$. By Lemma~\ref{lem:one-step-basic}(a) and Proposition~\ref{prop:avg-regular}, $|H|\le2da\delta_j^2\le2dc'_3$. Let $\varphi(t)=\ell(e^{\ii tZ_1}e^{\ii H}e^{-\ii tZ_2})-H$. For $|t|\le r_0/(|Z_1|+|Z_2|)$, with $r_0$ small depending on $\mathsf N$ (the logarithm is again a Baker--Campbell--Hausdorff series of $\mathfrak g^c$-elements, so $c_G$ does not enter), it is holomorphic and bounded by a constant $K_{\rm S}(\mathsf N)$, and $\varphi(0)=0$; this uses $2dc'_3\le r_0$. The Schwarz lemma gives $|\varphi(1)|\le(K_{\rm S}/r_0)(|Z_1|+|Z_2|)$. Average over $x$ for $E$.
\end{proof}

\begin{proposition}[Averaged backgrounds at complex references, (AB)]\label{prop:AB}
There are $\delta_{AB}>0$ and $c'_3\le c_3$, depending only on $d,L,G$ (through the constants of Part~IV, including $c_G$), such that the following holds for $a\le c'_3$, every real $U_r\in\mathfrak U_0(a)$, and every complex reference $U=e^{\ii\eta A'}U_r$ with $|A'|_\infty\le\delta_c<\delta_{AB}$. All bounds are uniform in $k$, in $\eta$ and in $j\le k$. Below, relative configurations, $\tau_X$, $\tilde{\mathfrak a}_j$, $\mathcal Q_j$, $l_j$, $q_j$, $r_j$, $\gamma_j$, $q'_j$, $\bQ_k$, $Q'_k$ at the reference $U$ are defined as in Section~\ref{sec:adapted}, with the complex backgrounds $W_j(U)$ in place of $W_j$.
\begin{enumerate}[label=(\roman*)]
\item \emph{(Backgrounds.)} $W_j(U)$ is defined for $j\le k$, and Lemmas~\ref{lem:AB-rep} and~\ref{lem:AB-bounds} hold. In particular, relative to $W_j(U_r)$, $W_j(U)$ is a complex perturbation of size at most $C_{\rm AB}\delta_c$, in every real gauge.
\item \emph{(One step.)} Lemma~\ref{lem:one-step-expansion} holds at $W=W_{j-1}(U)$, with $c_2$ replaced by $\frac12c_2$. Parts (a), (c), (d) are unchanged and (e) is dropped. Part (b) is replaced by
\begin{gather*}
|(q_jB)(c)|^2\le e^{2\kappa_{\rm AB}\delta_j\delta_c}L^{-d}\sum_xL^{-1}\sum_{b\in[x,x(c)]}|B(b)|^2,\\
|(r_jB)(c)|\le C(a\delta_j^2+\delta_j\delta_c)\max|B|,\qquad|\gamma_jB|\le de^{\kappa_{\rm AB}\delta_j\delta_c}\max|B| .
\end{gather*}
\item \emph{(Gauge-adapted averaging.)} Proposition~\ref{prop:Q-expansion} holds at the reference $U$, for complex $A$ with $|A|_\infty\le c'_3$. Parts (a), (b) and (d) hold in the same form, and ``real for real $A$'' is dropped. In (c) the kernel bounds become
\begin{gather*}
|K(c,b)|\le(1+Ca)\,e^{C_L\delta_c},\qquad|K_A(c,b)|\le Ce^{C_L\delta_c}|A|_\infty,\\
|d^2C_k(A)[h,h](c)|\le Ce^{C_L\delta_c}\|h\|^2_{S_c},
\end{gather*}
and the consequences of (c) (pointwise bounds by $\|h\|_{S_c}$, sup-norm bounds of the transposes) hold with their constants multiplied by $e^{C_L\delta_c}$.
\item \emph{(Linearized averages.)} $Q'_k(U)Q'_k(U)^T=I$, with the transpose taken with respect to $\mathrm b$. Proposition~\ref{prop:linearized-averages}(d), and (f) for real $u$, hold as stated. The tree transports satisfy $\|T_U(y,x)\|\le e^{\kappa_{\rm AB}\delta_c}$ and $\|T_U(y,x)-T_{U_r}(y,x)\|\le e^{\kappa_{\rm AB}\delta_c}-1$. Hence $|Q'_k(U)\lambda(y)|\le e^{\kappa_{\rm AB}\delta_c}\|\lambda\|_{B(y)}$, and $Q'_k(U)^T$ is bounded on sup norms by $e^{\kappa_{\rm AB}\delta_c}$. Proposition~\ref{prop:linearized-averages}(e) holds with $\eps$ replaced by $\eps+\delta_c$. More precisely, suppose that on a cube containing $B(c_-)\cup B(c_+)$ we have $U_r=e^{\ii\eta A_r}$ in a real gauge, with $|A_r|\le\eps$ and $a\le\eps\le c'_3$. Then, in this gauge, $|K_{\bQ_k(U)}(c,b)-K_{\bQ}(c,b)|+|K_{Q'_k(U)}(y,x)-K_{Q_0}(y,x)|\le C_L(\eps+\delta_c)$, and the same holds for the transposes.
\item \emph{(Closeness to the real reference.)} With the supports of Proposition~\ref{prop:linearized-averages}(d),
\begin{gather*}
|K_{\bQ_k(U)}(c,b)-K_{\bQ_k(U_r)}(c,b)|\le C_L\delta_c,\qquad |K_{Q'_k(U)}(y,x)-K_{Q'_k(U_r)}(y,x)|\le e^{\kappa_{\rm AB}\delta_c}-1,\\
\|R(W_k(U)(b))-R(W_k(U_r)(b))\|\le e^{2C_{\rm AB}\delta_c}-1 .
\end{gather*}
So $\bQ_k(U)-\bQ_k(U_r)$, $Q'_k(U)-Q'_k(U_r)$, $D_{W_k(U)}-D_{W_k(U_r)}$ and their transposes are local operators of size $C_L\delta_c$.
\item \emph{(Analyticity.)} $W_j(U)$, $l_j$, $\bQ_k(U)$, $Q'_k(U)$ and $D_{W_k(U)}$ are holomorphic in $A'$ on $\{|A'|_\infty<\delta_{AB}\}$. The nonlinear part $C_k(A)$ at the reference $U$ is holomorphic in $(A',A)$ jointly.
\end{enumerate}
\end{proposition}

\begin{proof}
(i) is Lemmas~\ref{lem:AB-rep} and~\ref{lem:AB-bounds}.

(ii) The proof of Lemma~\ref{lem:one-step-expansion} uses the reality of $W$ in four places; everything else is algebraic, or uses only analyticity and Cauchy estimates.
\begin{enumerate}[label=(\arabic*)]
\item \emph{Analyticity.} Relative transports are products of at most $3dL$ factors $e^{\ii\delta R(W(\gamma))B(b)}$, with $\gamma$ a contour of at most $3dL$ bonds of $\T^{(j-1)}$. By Lemma~\ref{lem:AB-bounds}(c), $\|R(W(\gamma))\|\le e^{\kappa_{\rm AB}\delta_j\delta_c}\le2$ instead of $1$. This is absorbed by $c_2\to\frac12c_2$. By Lemma~\ref{lem:AB-bounds}(d), the arguments of all logarithms are within $C(c_2+a\delta_j^2+\delta_j\delta_c)$ of $1$. The Cauchy estimates for $c_j$, $c^\tau_j$ are then unchanged.
\item \emph{The loops.} In the linear part, $|H_x|,|E|\le2da\delta_j^2$ is replaced by $2da\delta_j^2+C_{\rm loop}\delta_j\delta_c$, by Lemma~\ref{lem:AB-bounds}(d). The bounds $|\mathfrak f(\operatorname{ad}H_x)-1|\le C|H_x|$, $|\mathfrak e(\operatorname{ad}E)-1|\le C|E|$ and $|R(e^{\ii E})-R(h_x)|\le C(|E|+|H_x|)$ hold for complex $H_x,E$ of size at most $1$. Indeed $\mathfrak e$ is entire, $\mathfrak f$ is analytic for $\|\operatorname{ad}\|<2\pi$ (Lemma~\ref{lem:one-step-expansion}), and $\|\operatorname{ad}H\|\le2|H|\le2$. Moreover $|R(W(\Gamma^j_{c,x}))-R(W_j(c))|=\|(R(h_xe^{-\ii E})-1)R(W_j(c))\|\le C(|E|+|H_x|)$, since $\|R(W_j(c))\|\le2$.
\item \emph{The tube average.} In $q_jB(c)$ every value $B(b)$ is multiplied by a single transport along a contour of at most $(d+1)L$ bonds of $\T^{(j-1)}$, which has norm at most $e^{\kappa_{\rm AB}\delta_j\delta_c}$. The Cauchy--Schwarz inequality gives the bound on $q_j$; the bound on $\gamma_j$ is the same.
\item \emph{The remainder $r_j$.} The algebra of the proof is unchanged: with all factors $\mathfrak e,\mathfrak f,R(h_x),R(e^{\ii E})$ equal to $1$ and $R(W(\Gamma^j_{c,x}))=R(W_j(c))$, only $q_jB$ remains. So $r_jB$ is a sum of terms, each containing one of the differences in (2), multiplied by a contour sum bounded by $3dLe^{\kappa_{\rm AB}\delta_j\delta_c}\max|B|$, and by the prefactor $L^{-1}$. Hence $|r_jB|\le C(a\delta_j^2+\delta_j\delta_c)\max|B|$.
\end{enumerate}

(iii) \emph{(a), (b).} By (ii), $|l_jB|_\infty\le\big(1+C(a\delta_j^2+\delta_j\delta_c)\big)|B|_\infty$ and $|c_j(B)|_\infty\le C\delta_j|B|^2_\infty$. The induction of the proof of Proposition~\ref{prop:Q-expansion} gives, by \eqref{eq:AB-scales} and $\sum_i\delta_i^2\le\frac43$,
\[
|B_j|_\infty\le|A|_\infty\prod_{i\le j}\big(1+C(a\delta_i^2+\delta_i\delta_c)+2C\delta_i|A|_\infty\big)\le|A|_\infty e^{C(2a+2\delta_c+4|A|_\infty)}\le2|A|_\infty
\]
for $a,|A|_\infty\le c'_3$ and $\delta_c<\delta_{AB}$. The bounds on $C_j$ and the locality follow as there.

\emph{(c).} We run the proof of Proposition~\ref{prop:Q-expansion}(c) with nonnegative majorants. For a scale-$(m-1)$ field $v$ write $|v|$ for the function $b\mapsto|v(b)|$. In $q_mv(c')$ each value $v(b)$ is multiplied by one transport of norm at most $e^{\kappa_{\rm AB}\delta_m\delta_c}$, by (ii)(3), and the scalar weights are those at $W=1$. So $|q_mv|\le e^{\kappa_{\rm AB}\delta_m\delta_c}q^{(1)}_m|v|$ pointwise, where $q^{(1)}_m$ is the tube average at $W=1$. By (ii)(4), $|r_mv|\le r^{\rm maj}_m|v|$, where $r^{\rm maj}_m$ is a nonnegative kernel with the support of the real case and at most $C_L$ entries per row, each bounded by $C(a\delta_m^2+\delta_m\delta_c)$; it already accounts for the complex transports inside $r_m$. Hence
\[
|l_j\cdots l_{i+1}v|\le e^{\kappa_{\rm AB}\delta_c\sum_m\delta_m}\big(q^{(1)}_j+r^{\rm maj}_j\big)\cdots\big(q^{(1)}_{i+1}+r^{\rm maj}_{i+1}\big)|v| .
\]
The prefactor is at most $e^{2\kappa_{\rm AB}\delta_c}$, by \eqref{eq:AB-scales}. The product of majorants is the real case with unitary transports, except that the coefficients of the insertions are $C(a\delta_m^2+\delta_m\delta_c)$ instead of $Ca\delta_m^2$. The steps ``Transports'' and ``Insertions'' of that proof therefore apply verbatim, with $a\delta_m^2$ replaced by $a\delta_m^2+\delta_m\delta_c$. They show that the product is of type
\[
\prod_m\big(1+C_L(a\delta_m^2+\delta_m\delta_c)\big)\le\prod_m(1+C_La\delta_m^2)\prod_m(1+C_L\delta_m\delta_c)\le(1+Ca)\,e^{C_L\delta_c}.
\]
In the step ``First derivatives'', $dc_m(B)$ plays the role of $r_m$, with coefficients $C\delta_m|A|_\infty$ by the Cauchy estimates of (ii), and the same majorant argument applies. In the step ``Second derivatives'', the type bounds just proved replace those of the real case. So both steps hold with constants multiplied by $e^{C_L\delta_c}$. The consequences follow from the kernel bounds as before.

\emph{(d).} The values of $\mathsf w_k$ at the points of $\T$ are given by Theorem~\ref{thm:correspondence}(c), which is the recursion of Definition~\ref{def:adapted}, and by (ii); this gives the same bounds. At fine points, the extensions $\mathcal E^{(i)}$ transport along the trees, with norm at most $e^{\kappa_{\rm AB}\delta_c}$ instead of $1$.

(iv) For every $Y\in G^c$, $R(Y)^T=R(Y^{-1})$ with respect to $\mathrm b$, since $\operatorname{tr}(YXY^{-1}Z)=\operatorname{tr}(XY^{-1}ZY)$. Hence $q'^{\,T}_jg(x)=R(W_{j-1}(U)(\Gamma^j_{y,x}))^{-1}g(y)$ and $q'_jq'^{\,T}_j=I$, as in the proof of Proposition~\ref{prop:linearized-averages}(a). That proof uses only the volume factors, not unitarity. The Hermitian identity $Q'_kQ'^*_k=I$ need not hold for complex $U$; the $\mathrm b$-transpose identity always holds. Locality (d) holds because Lemma~\ref{lem:one-step-basic}(c) holds verbatim for complex configurations, and Lemma~\ref{lem:one-step-expansion}(a) holds by (ii). Covariance (f) follows from Lemma~\ref{lem:AB-cov} with real $g$.

The tree transport $T_U(y,x)$ is a product over $m<k$ of $R(W_m(U)(\gamma_m))$, with $\gamma_m=\Gamma^{m+1}_{x_{m+1},x_m}$ a contour of at most $d(L-1)$ bonds of $\T^{(m)}$. By Lemma~\ref{lem:AB-bounds}(c), $W_m(U)(\gamma_m)=e^{\ii Z_m}W_m(U_r)(\gamma_m)$ with $|Z_m|\le2dC_{\rm AB}\delta_{m+1}\delta_c$. By \eqref{eq:Ad-bound} and \eqref{eq:AB-scales}, $\|T_U\|\le e^{8dC_{\rm AB}\delta_c}\le e^{\kappa_{\rm AB}\delta_c}$. Telescoping the product gives $\|T_U-T_{U_r}\|\le e^{\kappa_{\rm AB}\delta_c}-1$.

For (e), compute in the given gauge, extended arbitrarily to $\Tf$; this is allowed by (f). By locality, $\tilde B_j$ on the relevant bonds depends only on $U_r,A'$ in the cube. By Lemma~\ref{lem:avg-local-gauge}, $W_j(U_r)=e^{\ii\delta_jA_{r,j}}$ with $|A_{r,j}|\le\eps+4da\le C\eps$. By Lemma~\ref{lem:AB-bounds}(a),(b) and Lemma~\ref{lem:log-estimates}(c), $W_j(U)=e^{\ii\delta_jA_{0,j}}$ with $|A_{0,j}|\le C(\eps+\delta_c)$. All transports along contours of at most $3dL$ bonds of $\T^{(j-1)}$ are then within $C(\eps+\delta_c)\delta_j$ of $1$. With (ii), $l_j-q_j(1)$ is a combination of at most $C_L$ values with coefficients $C(\eps+\delta_c)\delta_j$, using $a\delta_j^2\le\eps\delta_j$. The telescoping of the proof of Proposition~\ref{prop:linearized-averages}(e), with the type bounds of (iii), gives the kernel bound $\sum_jC_L(\eps+\delta_c)\delta_je^{C_L\delta_c}\le C_L(\eps+\delta_c)$. The same argument applies to $Q'_k$.

(v) The bound for $Q'_k$ is the tree bound in (iv). For the unit-lattice derivative, $W_k(U)(b)=e^{\ii\tilde B_k(b)}W_k(U_r)(b)$ with $|\tilde B_k|\le C_{\rm AB}\delta_c$, and \eqref{eq:Ad-bound} applies. For $\bQ_k$ telescope:
\[
\bQ_k(U)-\bQ_k(U_r)=\sum_{j\le k}l_k(U)\cdots l_{j+1}(U)\,\big(l_j(U)-l_j(U_r)\big)\,l_{j-1}(U_r)\cdots l_1(U_r).
\]
By the formula for $l_j$ in the proof of Lemma~\ref{lem:one-step-expansion}, $l_jB(c)$ is a sum of at most $C_L$ terms. Each is a product of a bounded number of factors from $\{R(W(\gamma)),R(e^{\ii E}),R(h_x),\mathfrak e(\operatorname{ad}E),\mathfrak f(\operatorname{ad}H_x)\}$ applied to one value $B(b)$, with $\gamma$ a contour of at most $3dL$ bonds. Each factor is bounded by $2$. By Lemma~\ref{lem:AB-bounds}(c),(d), each factor changes by at most $C\delta_j\delta_c$ from $U_r$ to $U$. So $l_j(U)-l_j(U_r)$ is a combination of at most $C_L$ values with coefficients $C_L\delta_j\delta_c$. With the type bounds of (iii), at $U$ and at $U_r$, the kernel of the sum is bounded by $\sum_jC_L\delta_j\delta_c(1+Ca)e^{C_L\delta_c}\le C_L\delta_c$. The supports are those of Proposition~\ref{prop:linearized-averages}(d). Transposes follow by summing the kernel bounds.

(vi) By Lemma~\ref{lem:AB-bounds}(b), $W_j(U)$ is holomorphic in $A'$. On the domain of (ii), $\delta_j^{-1}\ell(\tilde{\mathfrak a}_j(e^{\ii\delta_{j-1}B})(c))$ is jointly holomorphic in $B$ and in the finitely many bond variables of $W_{j-1}(U)$ on which it depends. Its derivative at $B=0$, which is $l_j$, is therefore holomorphic in $A'$. The rest follows by composition, uniformly in $k$ by the bounds of (iii)--(v).
\end{proof}

\begin{remark}[Why the naive bound fails]\label{rem:AB-outline}
The linear part in Lemma~\ref{lem:AB-bounds}(a) explains why a per-scale bound fails. The derivative at $A'=0$ of the plain iterated average is $\bQ_j-D_{W^r_j}\boldsymbol\Gamma_j$: the gauge-adapted linearization plus a coarse pure-gauge term. The plain one-step relative average has a sup-norm loss of up to a factor $2d+1$, which comes entirely from that pure-gauge term. The term does not accumulate, since $\boldsymbol\Gamma_j$ is a sum of $\delta_{i+1}$-weighted contour averages.
\end{remark}

\begin{remark}[Scope of Proposition~\ref{prop:AB}]\label{rem:AB-remains}
\begin{enumerate}[label=(\arabic*)]
\item \emph{Not proved, and not needed below.}
  \begin{itemize}
  \item A curvature bound for $W_j(U)$ of order $a+\delta_c$ in field-strength units for $j<k$, i.e.\ a complex analogue of Lemma~\ref{lem:one-step-regularity} and Proposition~\ref{prop:avg-regular}. We only have $|H_x(U)|\le2da\delta_j^2+C_{\rm loop}\delta_j\delta_c$. This is weaker by $\delta_j^{-1}$ in the $\delta_c$-term, but summable over scales. At the unit scale the two agree: $|F_{W_k(U)}|\le2a+C\delta_c$, from Lemma~\ref{lem:AB-bounds}(c) with $n=4$ and the Schwarz-lemma argument of (d).
  \item Derivative or H\"older bounds of $\tilde B_j$ for $0<j<k$.
  \item Theorem~\ref{thm:correspondence}(a),(b),(e),(f) and Lemma~\ref{lem:axial}(b),(c) at complex references. They are stated for real relative configurations; only Theorem~\ref{thm:correspondence}(c),(d) and Lemma~\ref{lem:axial}(a) hold for complex ones.
  \end{itemize}
  Lemma~\ref{lem:one-step-expansion}, Proposition~\ref{prop:Q-expansion} and Proposition~\ref{prop:linearized-averages} use the backgrounds only through transports along contours of bounded length. If a later argument needs any of the items above, it must be supplied separately.
\item \emph{Level of detail.} Proposition~\ref{prop:AB}(ii)--(v) is a line-by-line modification of the proofs in Section~\ref{sec:adapted}. It inherits their level of detail, in particular the type bookkeeping of Proposition~\ref{prop:Q-expansion}(c). The new inputs, Lemmas~\ref{lem:AB-cov}--\ref{lem:AB-bounds}, are self-contained given Proposition~\ref{prop:Q-expansion} at real references.
\item \emph{Downstream.}
  \begin{itemize}
  \item The bounds $|\bQ A(c)|\le C\|A\|_{S_c}$ and $|Q'\lambda(y)|\le C\|\lambda\|_{B(y)}$, and the sup-norm bounds of $\bQ^T,Q'^T$, used in Theorem~\ref{thm:ID} and Part~IX, follow from (iii) and (iv) with constants $e^{C_L\delta_c}\le2$.
  \item The exact identities ($E=0$, Proposition~\ref{prop:linearized-averages}(b), and those listed in step~(iii) of the proof of Proposition~\ref{prop:complex-ext}) are not extended on all of $\{|A'|_\infty<\delta_{AB}\}$: there only the sup norm of $A'$ is controlled, so the real points need not be regular backgrounds. They hold on the ball $\mathcal B$ of that proof, whose real points are regular, by (vi) and the identity theorem.
  \end{itemize}
\end{enumerate}
\end{remark}

\subsection{Extension of the results of Part~IX}\label{sec:complex-ext}

\begin{proposition}[Extension to complex references]\label{prop:complex-ext}
Let $\eps_{\max}$ be half the smallest of the real-reference thresholds on $\eps$ that are used in the proof below, each computed with the multiplicative constants doubled. There is $\delta_P>0$, depending on $L,M,\kappa,\rho,\alpha$, on the loss parameter $\delta$ of Part~IX, on $\eps_{\max}$ and on $G$, such that the following holds for $\delta_c<\delta_P$, at complex references of size $\delta_c$ whose real part is $(M,\eps)$-regular with $\eps\le\eps_{\max}$. The statements hold with $\eps$ replaced by $\eps+C\delta_c$. The $\Rea$-coercivity constants and the decay rates do not depend on $\delta_c$. The coercivity constants are those of Proposition~\ref{prop:reg-complex}, i.e.\ $\frac c2$ with $c$ that of the real reference (Lemma~\ref{lem:CT-stable}), and at level~(L5) $\frac18c_L$ and $\frac1{16}c_L$ (Lemma~\ref{lem:coercive-general}, Proposition~\ref{prop:local-fluct}). The decay rates are those that the real-reference proofs give when every input decay bound is taken at the fixed rate of the differences, a quarter of the minimum of the perturbation rate and the two chosen inverse decay rates in Lemma~\ref{lem:perturbed-inverse} (the step ``The constants'' of the proof). The other constants, computed at the real reference with these coercivity constants and rates, change by at most a factor $1+C_{\rm cx}\delta_c\le2$, where $C_{\rm cx}$ depends on $L,M,\kappa,\delta,\rho,\alpha$ and $G$, but not on $\delta_c$, and $C_{\rm cx}\delta_P\le1$:
\begin{enumerate}[label=(\alph*)]
\item Theorem~\ref{thm:ID}(a)--(c), Corollary~\ref{cor:lp-dirichlet}, and Lemmas~\ref{lem:int-factors}, \ref{lem:lift}, \ref{lem:N-tier};
\item Lemma~\ref{lem:trace}, Lemma~\ref{lem:cut-term}, and Theorems~\ref{thm:tier-walk} and~\ref{thm:HC-pointwise};
\item Proposition~\ref{prop:action}(a)--(c). Here $F^s$ and $J=D^TF^s$ are the holomorphic continuations, with $P_{\mathfrak g}$ replaced by its complex-linear extension $P_{\mathfrak g^c}$ and the summand continued without taking real parts. $|F^s|\le C(\eps+\delta_c)$, while $J$ need not be bounded uniformly in $\eta$;
\item Theorem~\ref{thm:weak-bg}(a), (d), and (b) without the realness statement. The fixed point $A_*(s;B)$ is analytic in $A'$, jointly with $(s,B)$.
\end{enumerate}
Not extended, and not needed at complex references: Theorem~\ref{thm:ID}(d), the identity of Lemma~\ref{lem:div-curvature} (its bound on $W_0$ extends; proof of Proposition~\ref{prop:dirichlet-hessian}(c)), Lemma~\ref{lem:critical-current} and Theorem~\ref{thm:weak-bg}(c). The ``real for real arguments'' statements of Propositions~\ref{prop:Q-expansion} and~\ref{prop:action} and of Lemma~\ref{lem:one-step-expansion}(e) are dropped at complex references as well.
\end{proposition}

\begin{proof}
We go through the places where reality enters.
\begin{enumerate}[label=(\roman*)]
\item \emph{The $L^2$ theory.}
  \begin{itemize}
  \item Proposition~\ref{prop:reg-complex} gives coercivity (as $\Rea$-coercivity), decay and comparison of domains for the Dirichlet operators of Parts~III and~VIII. It was proved for the averaging of Section~\ref{sec:background}; for the gauge-adapted averaging of Part~IV it needs property~(v) of Remark~\ref{rem:general-averaging}, i.e.\ Proposition~\ref{prop:AB}(iv)--(vi) and step~(iii) below. The averaged operators $\boldsymbol\Gamma^X$, $\bH^X$, $\bC^X$ follow by compositions and perturbed inverses.
  \item Part~IX allows complex perturbations from the outset (Section~\ref{sec:weakenable}, Setting).
  \end{itemize}
\item \emph{The local gauge in the interior estimates.} Instead of a real gauge in which $U$ satisfies (B), use the real gauge of $U_r$. There $U=e^{\ii\eta A_0}$ with complex $A_0$ satisfying (B) with parameter $C(\eps+\delta_c)$.
  \begin{itemize}
  \item Lemmas~\ref{lem:cov-expansion} and~\ref{lem:curvature} allow $\mathfrak g^c$-valued $A$.
  \item The absorption, iteration and $L^p$ arguments use only sup norms and quotients, no maximum principle.
  \item The flat estimates (Theorems~\ref{thm:local-regularity}, \ref{thm:lp-local}, Lemmas~\ref{lem:rescaled}, \ref{lem:lp-kernels}) apply to complex-valued fields componentwise.
  \end{itemize}
\item \emph{Algebraic and analytic identities.} Fix $U_r$ and consider $A'$ in the open ball $\mathcal B=\{\text{size}(A')<\delta_P\}$. The size is a norm on the complex fields $A'$, so $\mathcal B$ is convex. Real $A'$ give real configurations $U=e^{\ii\eta A'}U_r$ that are $(M,\eps+C\delta_P)$-regular, with curvature at most $C(\eps+\delta_P)$; this needs $\eps_{\max}+C\delta_P\le\min(\eps_0,c'_3)$. At these real points the operators below are those of Parts~IV--IX at the real background $U$: Proposition~\ref{prop:AB} defines all objects at the reference $U$ through $W_j(U)=U^{(j)}$ by the formulas of Section~\ref{sec:adapted}, and inverses are unique. So every identity proved in the note for real regular backgrounds of small curvature holds at the real points of $\mathcal B$. These form a totally real subspace of maximal dimension of the convex set $\mathcal B$. So an identity between two holomorphic functions on $\mathcal B$ that holds at the real points holds on all of $\mathcal B$, by the identity theorem. This is used in two stages.
  \begin{itemize}
  \item[(iii-a)] Identities between objects of Proposition~\ref{prop:AB} and of the action. These are holomorphic on $\mathcal B$ by Proposition~\ref{prop:AB}(vi), and since the action is a polynomial in the entries of $U$ and $U^{-1}$. They are the exact covariance $E=0$ (Proposition~\ref{prop:linearized-averages}(b)); $J=D^TF^s$ (Lemma~\ref{lem:current}(a)); and $\ip{\rho^{(1)}(A)}{F^s}=0$ in the proof of Proposition~\ref{prop:action}. After this stage the pair $(Q'_k,\bQ_k)$ has property~(v) of Remark~\ref{rem:general-averaging} on $\mathcal B$, and step~(i) applies. Stage~(iii-a) holds with the same proof on every such ball whose real points satisfy $\eps+C\cdot(\text{radius})\le\min(\eps_0,c'_3)$; in particular on the ball of radius $2\delta_*$ of the proof of Proposition~\ref{prop:dirichlet-hessian}(c).
  \item[(iii-b)] Identities that involve inverses. After step~(i), the operators of Parts~III, VIII and~IX are holomorphic on $\mathcal B$, as compositions and inverses of holomorphic families with uniformly bounded inverses, and uniformly convergent series of such. The identities are those of Lemma~\ref{lem:domain-identities} used in Lemma~\ref{lem:local-system}; $\bG DR=DG'R$ (Lemma~\ref{lem:Api-identities}(a)) and \eqref{eq:C-identity}; the $\mathrm b$-symmetry $(\bG^Z)^T=\bG^Z$ used in Lemma~\ref{lem:cut-term}; and the identity of Proposition~\ref{prop:one-loop-polymer}(a). For the last, $\log\det\Delta_k$ means $\sum_b\ell_k(b)$, with the integral of Proposition~\ref{prop:det-densities}(c). On $\mathcal B$, $\Rea\Delta_k\ge\frac18c_L$ (Lemma~\ref{lem:coercive-general}; it needs $|J_{U_r}|\le\eps$, which holds at every reference where Proposition~\ref{prop:one-loop-polymer} is applied, Section~\ref{sec:step-quantities}, Setting; its conditions $\eps\le\eps_{\rm cc}$ and those of the proof of Proposition~\ref{prop:dirichlet-hessian}(c) hold since $\eps\le\eps_{\max}$), so the spectrum of $\Delta_k$ lies in $\{\Rea z>0\}$. There the integral is the principal logarithm: it is holomorphic on $\mathcal B$, real at the real points, and its exponential is $\det\Delta_k$.
  \end{itemize}
  No connectedness of the set of all complex references is needed. That set need not be connected: if $G$ has a $U(1)$ factor, real regular backgrounds with different fluxes exist.
\item \emph{Unitarity of transports and the averaging.} This is Proposition~\ref{prop:AB}.
\item \emph{The action.} The proof of Proposition~\ref{prop:action} uses a gauge with $|A_0|,|\partial A_0|\le\eps$, $|R(H_p)-1|\le C\eps\eta^2$ and $|F^s|\le C\eps$. These hold with $C(\eps+\delta_c)$, using the holomorphic summand and $P_{\mathfrak g^c}$. The current $J$ is not bounded uniformly in $\eta$. In the weakened background map it enters only as $(D\bH(s))^T\mathsf F$ (Section~\ref{sec:weak-background}).
\item \emph{The trace inequality.} Lemma~\ref{lem:trace} is stated for complex perturbations.
\item \emph{Analyticity in $A'$.} Every input is holomorphic in $A'\in\mathcal B$, by (i)--(v), and the contraction of Theorem~\ref{thm:weak-bg} is uniform in $A'$. The factorization of Theorem~\ref{thm:weak-bg}(d), including its statement on the dependence on $s$, follows from the same proof: Lemma~\ref{lem:weak-basic}(b) allows complex backgrounds, the averaging operators are local at complex references (Proposition~\ref{prop:AB}(iii),(iv)), and the nonlinear part is local by (c).
\item \emph{The constants.} The $\Rea$-coercivity constants and the decay rates at complex references do not depend on $\delta_c$. \emph{The fixed rates.} For the operators of Parts~II, III and~VIII, Lemma~\ref{lem:CT-stable} gives the hypotheses of Lemma~\ref{lem:CT-form} with the constants $c/2$, $C_1+1$ and $\min(\alpha_1,\gamma/2,1)$, where $c,C_1,\alpha_1$ are those at the real reference and $\gamma$ is the rate of the perturbation. So these operators decay at the rate $\gamma_{\rm c}=\min(\alpha_1,\gamma/2,\frac13,c/(4(C_1+1)))$. By Lemma~\ref{lem:perturbed-inverse}, their differences from the operators at $U_r$ have decay $(C\delta,\gamma_{\rm c}/4)$, with $\delta\le C_{\rm cmp}\delta_c$. At level~(L5) the lower bounds are $\frac18c_L$ and $\frac1{16}c_L$, and the differences again decay at a quarter of the minimum of the perturbation rate and the two chosen inverse decay rates. Every input decay bound below, at real and at complex references, is taken at these quarter rates. At the rate $\gamma_{\rm c}$ itself the constants at a complex reference are only bounded by about twice the real ones.

  \emph{The other constants.} Order the levels of Sections~\ref{sec:nested-walks} and~\ref{sec:tiers}, and the constants of Section~\ref{sec:weak-background}, as in Section~\ref{sec:nested-walks}: $\ell=1,\dots,N$. Each multiplicative constant $K$ of level $\ell$ is $F_K(x_\ell)$, where $x_\ell$ collects the primitive inputs and the constants of the earlier levels, and $F_K$ is a nonnegative combination of products of nonnegative, nondecreasing, positively $1$-homogeneous functions of $x_\ell$. These include sums, products, maxima, Neumann and walk series with nonnegative coefficients, and integrals of such, like the resolvent integral of Corollary~\ref{cor:weak-covariance}. So $t\mapsto F_K((1+t)x_\ell)$ is convex. Put $\sigma_{\rm cx}=2^{1/3}-1$, so that $(1+\sigma_{\rm cx})^3\le2$. For $x'\le(1+t)x$ and $0\le t\le\sigma_{\rm cx}$, convexity gives
  \[
  F_K(x')\le F_K(x)\Big(1+\frac t{\sigma_{\rm cx}}\Big(\frac{F_K((1+\sigma_{\rm cx})x)}{F_K(x)}-1\Big)\Big).
  \]
  At a complex reference the primitive inputs change by at most a factor $1+C_{\rm in}\delta_c$: the averaging kernels by $e^{C_L\delta_c}\le1+2C_L\delta_c$ (Proposition~\ref{prop:AB}(iii)--(v)); the operators of Parts~II, III and~VIII by the differences above; and $W_2^X$ by $C_{\rm cmp}C_W\delta_c$ (proof of Proposition~\ref{prop:dirichlet-hessian}(c)). Each input bound $K$ is taken free of $\eps$ (with $\eps\le1$; e.g.\ $\|W_2^X\|\le C$), so an additive change $C\delta_c$ is a factor $1+(C/K)\delta_c$; for inverses, Lemma~\ref{lem:perturbed-inverse} gives the change relative to the bound, also at level~(L5), where $C_G<1$. Put $C^{(\ell)}=(C^{(\ell)}_{\rm in}/\sigma_{\rm cx})\max\big(1,\max_K(F_K((1+\sigma_{\rm cx})x_\ell)/F_K(x_\ell)-1)\big)$, the maximum over the constants of level $\ell$, with $C^{(\ell)}_{\rm in}=\max(C_{\rm in},C^{(1)},\dots,C^{(\ell-1)})$. Let $C_{\rm cx}=\max_\ell C^{(\ell)}$. For $\delta_c\le1/C_{\rm cx}$ one has $C^{(\ell)}_{\rm in}\delta_c\le\sigma_{\rm cx}$ for every $\ell$. So, by induction over $\ell$, every input of every level changes by at most a factor $1+\sigma_{\rm cx}$, and every constant by at most $1+C_{\rm cx}\delta_c\le2$. The values $F_K((1+\sigma_{\rm cx})x_\ell)$ are finite. Indeed, the step constants $\theta$ of Lemma~\ref{lem:step-small} and $\theta'$ of Theorem~\ref{thm:tier-walk} are sums of products of at most three inputs of their level, so their values at $(1+\sigma_{\rm cx})x_\ell$ are at most twice those at $x_\ell$. Moreover $M_1$ is computed with the inputs doubled, and $M_1'$ with $\theta'$ replaced by $2\theta'$. So for $M\ge M_1,M_1'$ the walks converge at $(1+\sigma_{\rm cx})x_\ell$, with ratio at most $\frac12$. For the same reason $2\theta'$ bounds $\theta'$ at complex references, and the inputs of $M_1$ at most double there. $C_{\rm cx}$ is built from real-reference constants that hold for all $\eps$ below thresholds fixed before $\delta_P$, such as $\eps_0$ and $c_5$. So it depends on $L$, $M$, $\kappa$, $\delta$, $\rho$, $\alpha$ and $G$, but not on $\delta_P$. The thresholds $M_1$ and $M_1'$, fixed in this way before $M$, serve at every complex reference with $\delta_c<\delta_P$, and $M$ is chosen before $\delta_P$ without circularity. The proofs of Theorems~\ref{thm:tier-walk} and~\ref{thm:HC-pointwise} then apply verbatim.
\item \emph{Not extended.}
  \begin{itemize}
  \item The identity of Lemma~\ref{lem:div-curvature} differentiates real gauge transformations and needs $|J_U|\le\eps$. It enters only Theorem~\ref{thm:ID}(d), which is not used after Section~\ref{sec:weak-background}.
  \item Lemma~\ref{lem:critical-current} and Theorem~\ref{thm:weak-bg}(c) concern real critical backgrounds, where they are used.
  \end{itemize}
\end{enumerate}
\emph{Smallness.} $\delta_P$ must satisfy the following:
\begin{itemize}
\item $C_{\rm AB}\delta_P\le\min(c'_L,cL^{-d-1})$, for the elimination of Lemma~\ref{lem:eliminate} at a complex base (Corollary~\ref{cor:covariance}, with $|\mathfrak r_c|\le\frac12$) and for the coercivity of Lemma~\ref{lem:coercive-general}, used in Section~\ref{sec:step-quantities}; and $\delta_P\le c'_L$. Since $C_{\rm cmp}c'_L\le\delta_*$, with $C_{\rm cmp}=4$ (first item of Section~\ref{sec:complex-references}) and $\delta_*$ of the proof of Proposition~\ref{prop:dirichlet-hessian}(c), the perturbations of size $<C_{\rm cmp}\delta_P$, and the ball of twice that size used there for the Cauchy estimate, lie within the radii fixed with $\delta_*$. These are the radii of Proposition~\ref{prop:reg-complex} for $W=0$ and for $W_2$, $W_2^X$, of Proposition~\ref{prop:AB}, of Proposition~\ref{prop:action}, of Proposition~\ref{prop:dirichlet-averaged}(e), and of $E=0$ (step~(iii));
\item $\eps_{\max}+C\delta_P\le\min(\eps_0,c'_3,\eps_L,\eps_{\rm ID},\eps_{\rm cc},\frac1{20})$, with $\eps_L$ the threshold of Proposition~\ref{prop:system-interior}, $\eps_{\rm ID}$ that of Theorem~\ref{thm:ID} and $\eps_{\rm cc}$ that of the first item of Section~\ref{sec:complex-references}; moreover $\eps_{\max}+C\delta_P$ is below the other thresholds of Sections~\ref{sec:interior} and~\ref{sec:interior-lp};
\item $C(\eps_{\max}+\delta_P)\le c_5$ of Proposition~\ref{prop:action};
\item $C_P(\eps_{\max}+C\delta_P)$ small in Theorem~\ref{thm:weak-bg}, with $C_P$ computed with the multiplicative constants doubled;
\item at complex references the condition $(1+K_H)r_P\le c_3$ in the proof of Lemma~\ref{lem:Ds} becomes $(1+2K_H)r_P\le c'_3$, with $r_P$ computed with the multiplicative constants doubled;
\item $C_{\rm cx}\delta_P\le1$, for the factor $1+C_{\rm cx}\delta_c\le2$.
\end{itemize}
At complex references the constant $K_H$ is at most $2K_H$, by the step ``The constants'' of the proof. Later sections use $2K_H$ in place of $K_H$, and the radii $R_P$, $r_P$ and the constants $C_P$, $\eps_P$, $\varrho_P$ of Theorem~\ref{thm:weak-bg} computed with the multiplicative constants doubled.
All of these are conditions on $\delta_P$ after $L,M,\kappa,\rho,\alpha$ and $\eps_{\max}$ are fixed.
\end{proof}

\begin{remark}[Use in the renormalization group]
In the activity domain D3$'''$ the complex perturbations have size at most $Ch_k^{1-2\epsilon}$, which tends to $0$ as $g_*\to0$. So, once $g_*$ is small, $\delta_c<\delta_P$, and the real parts are $(M,\eps)$-regular with $\eps\le\eps_{\max}$: by Corollary~\ref{cor:regularity}, $\eps\le B_5h_k^{1-\epsilon/2}$, with $B_5$ fixed before $g_*$. This is consistent with the order of the parameters in Theorem~\ref{thm:one-step}. In the units of step $k$ the budgets of the next step shrink further, which only helps. By Propositions~\ref{prop:AB} and~\ref{prop:complex-ext}, the new activities produced by the step are analytic functions of the complex background, as the norm requires.
\end{remark}
